\documentclass[11pt,a4paper]{article}
\pdfinfoomitdate=1
\pdftrailerid{}
\pdfsuppressptexinfo=15
\usepackage[T1]{fontenc}
\usepackage{lmodern,amsmath,amssymb,amsthm}
\usepackage[margin=27mm]{geometry}
\usepackage{booktabs}
\usepackage[hidelinks]{hyperref}
\emergencystretch=1em
\hypersetup{pdftitle={Ringmin Proposition 7: standalone lower-bound supplement},pdfauthor={Maurizio Falconi}}
\makeatletter
\renewcommand*\l@section{\@dottedtocline{1}{0em}{2.5em}}
\renewcommand*\l@subsection{\@dottedtocline{2}{2.5em}{3em}}
\makeatother
\newtheorem{theorem}{Theorem}[section]
\newtheorem{lemma}[theorem]{Lemma}
\newtheorem{proposition}[theorem]{Proposition}
\newtheorem{corollary}[theorem]{Corollary}
\theoremstyle{remark}\newtheorem{remark}[theorem]{Remark}
\newcommand{\Rstar}{R^*}
\newcommand{\Cterm}{C_{\rm term}}
\newcommand{\etaw}{\eta_{\rm width}}
\newcommand{\Rc}{R_{\rm chain}}
\newcommand{\ind}{\mathbf 1}
\newcommand{\dd}{\,\mathrm d}
\newcommand{\pos}[1]{\left(#1\right)_+}
\DeclareMathOperator{\supp}{supp}
\title{A standalone proof of the lower endpoint in Proposition 7\\
of the Ringmin asymptotic sequel}
\author{Maurizio Falconi}
\date{23 October 2026\\Journal supplement for independent review}
\begin{document}
\maketitle
\begin{abstract}
We give a self-contained proof of the explicit lower bound
$C_*\ge\Cterm+\etaw$ for the asymptotic central-circle constant of
Ringmin. Three nested terminal deletions of one cyclic order share one
quadratic energy. An explicit dual potential, signed deletion estimate,
crossing lemma, scalar minimax and global cubic certificate yield the bound.
All constants and directed rational checks are specified here and in two
small included arithmetic scripts. In particular,
\[
0.14056946887766098063257<\Cterm+\etaw
<0.14056946887766098063392.
\]
The coefficient is a lower bound, not an identified geometric optimum.
This supplement reconstructs Proposition 7 of arXiv:2609.13630 and does
not revise that manuscript or assert external scientific acceptance.
\end{abstract}
{\small\tableofcontents}
\clearpage

\section{Statement, notation and proof dependencies}\label{sec:statement}
The Ringmin problem places circles of radii $1,\ldots,n$ externally
tangent to a central circle of radius $R$, requiring the interiors of
every pair of outer circles to be disjoint. Write $\Rstar(n)$ for the
infimum of such feasible $R$. For a cyclic order $\omega$ of at least
three distinct positive radii, $\Rc(\omega)$ is the adjacent angular
closure root defined in Section~\ref{sec:geometry}. It is only a necessary
bound for a full feasible placement of that order.

Let $\tau$ be the unique solution of $\cos\tau=\tau$ in $(0,1)$ and put
\begin{equation}\label{eq:parameters}
q=\frac{1-\sin\tau}{1+\sin\tau},\qquad
\Cterm=\frac{\tau(1+q)}{2\pi},\qquad
(\beta_1,\beta_2,\beta_3)=
\left(\frac{2091}{10000},\frac{10907}{50000},\frac{23}{100}\right).
\end{equation}
Define
\begin{align}
D_i&=\int_q^{\beta_i}\left[1+q-x-2\sqrt{x(1+q-x)}\right]\dd x,
\label{eq:Di}\\
a&=\beta_2-\beta_1=\frac{113}{12500},\qquad
b=\beta_3-\beta_2=\frac{593}{50000},\nonumber\\
\mathcal P&=\{h\in[0,\infty)^3:h_1+h_2\le a,\ h_2+h_3\le b\},\nonumber\\
F(h)&=\sum_{i=1}^3h_iD_i-8\sum_{i=1}^3h_i^3,\qquad
\mathcal L(h)=16+432\sum_{i=1}^3h_i,\label{eq:objective}\\
\eta(h)&=\frac{\pos{F(h)}}{\pi\mathcal L(h)},\qquad
\etaw=\max_{h\in\mathcal P}\eta(h),\quad r_+=\max\{r,0\}.\nonumber
\end{align}

\begin{theorem}\label{thm:main}
The limit $C_*=\lim_{n\to\infty}\Rstar(n)/n^2$ exists. The maximum
in \eqref{eq:objective} is positive and uniquely attained at
$u=(a-z_*,z_*,b-z_*)$, where $z_*$ is the unique root on $[0,a]$ of
the polynomial $G$ defined in \eqref{eq:G}. Moreover
\begin{equation}\label{eq:mainbound}
\begin{gathered}
C_*\ge\Cterm+\etaw,\\
\frac{14056946887766098063257}{10^{23}}
<\Cterm+\etaw<
\frac{14056946887766098063392}{10^{23}}.
\end{gathered}
\end{equation}
\end{theorem}

Every terminating decimal in this document denotes an exact rational.
We use $\alpha=k/n$ for the finite grid's lower endpoint; $a,b$ are
reserved for the two width constraints. The root $z_*$ belongs solely
to this lower bound and has no connection to an upper-bound parameter.
The proof is organized as follows.

\begin{center}\small
\begin{tabular}{@{}cll@{}}\toprule
Input & Complete argument & Dependency\\\midrule
L1 & Angular geometry and deletion (\S\ref{sec:geometry}) & definitions\\
L2 & Exact terminal parameters (\S\ref{sec:terminal}) & elementary calculus\\
L3 & Dual and common energy (\S\ref{sec:dual}) & L2 grid domain\\
L4 & Signed deletion and floors (\S\ref{sec:deletion}) & L3\\
L5 & Shared crossings (\S\ref{sec:crossings}) & common measure\\
L6 & Finite minimax (\S\ref{sec:minimax}) & L1, L3--L5\\
L7 & Width optimum and enclosure (\S\ref{sec:width}, \S\ref{sec:arithmetic}) & L2\\
L8 & Ordered limits (\S\ref{sec:limit}) & L6, L7; limit existence\\\bottomrule
\end{tabular}
\end{center}
The final limit-existence argument is included in Section~\ref{sec:existence}.
The references identify provenance; no external proof or operational log is
needed to follow the arguments below.

\section{Angular geometry and deletion (L1)}\label{sec:geometry}
For radii $u,v>0$ the angle required by non-overlap is
\begin{equation}\label{eq:kernel}
\theta_R(u,v)=2\arcsin\frac{\sqrt{uv}}{\sqrt{(R+u)(R+v)}}
=2\arctan\frac{\sqrt{uv}}{\sqrt{R(R+u+v)}}.
\end{equation}
Indeed the centers lie at distances $R+u,R+v$ from the origin. The cosine
law says that their distance is at least $u+v$ exactly when the smaller
angular separation is at least \eqref{eq:kernel}. The second identity
uses $1-uv/((R+u)(R+v))=R(R+u+v)/((R+u)(R+v))$.
Both directed angular arcs must therefore be at least $\theta_R(u,v)$.

\begin{lemma}[Uniform root sandwich]\label{lem:sandwich}
For a cycle $\omega$ of at least three radii in $\{1,\ldots,n\}$, put
$S(\omega)=\sum_{\{u,v\}\in\omega}\sqrt{uv}$, counting every cyclic edge
once, including the closing edge. Then the equation
$\sum_{\omega}\theta_R(u,v)=2\pi$ has a unique positive root and
\begin{equation}\label{eq:sandwich}
\frac{S(\omega)}\pi-n\le\Rc(\omega)\le\frac{S(\omega)}\pi.
\end{equation}
\end{lemma}
\begin{proof}
Each angle is continuous and strictly decreasing from $\pi$ to zero as
$R$ increases from zero to infinity. The sum starts above $2\pi$.
For $1\le u,v\le n$, $\arcsin t\ge t$, $\arctan t\le t$ and
\eqref{eq:kernel} give
\[
\frac{2\sqrt{uv}}{R+n}\le\theta_R(u,v)\le\frac{2\sqrt{uv}}R.
\]
At the root these give $S/(R+n)\le\pi\le S/R$, which is precisely
\eqref{eq:sandwich}. No scale assumption or parity restriction is used.
\end{proof}

\begin{lemma}[Deletion from full geometry]\label{lem:deletiongeometry}
Delete any outer circles from a full feasible configuration, retaining
at least three circles with their original radii. If $\omega'$ is the
cyclic order of the survivors, then $\Rc(\omega')\le R$.
\end{lemma}
\begin{proof}
Every surviving pair remains non-overlapping. If $g$ is the directed gap
between successive surviving centers, then
$\min\{g,2\pi-g\}\ge\theta_R(u,v)$, hence $g\ge\theta_R(u,v)$.
All these gaps, including closure, sum to $2\pi$. Thus the chain sum
at $R$ is at most $2\pi$; strict decrease gives the assertion.
This proof uses full pair feasibility before deletion. It never reconstructs
a full placement from a chain root.
\end{proof}

\section{Exact terminal parameters (L2)}\label{sec:terminal}
\begin{lemma}\label{lem:terminal}
The parameters in \eqref{eq:parameters} are uniquely defined, with
$3/17<q<1/5<\beta_1$. Furthermore
\begin{equation}\label{eq:J}
J:=\int_q^1\sqrt{x(1+q-x)}\dd x=\pi\Cterm.
\end{equation}
Among $0<r<1$ the expression
$\pi^{-1}\int_r^1\sqrt{x(1+r-x)}\dd x$ has a unique global maximum
at $r=q$.
\end{lemma}
\begin{proof}
The function $\cos t-t$ is continuous, positive at zero, negative at one
and strictly decreasing on $(0,1)$, since its derivative is $-1-\sin t$.
This proves existence and uniqueness of $\tau$. The rational gates in
Section~\ref{sec:arithmetic} prove $3/17<q<1/5$ without assuming rounded
approximations.

For arbitrary $r\in(0,1)$ put $v=(1-r)/(1+r)$. Substitution
$y=(2x-1-r)/(1+r)$ gives
\[
\frac1\pi\int_r^1\sqrt{x(1+r-x)}\dd x
=\frac{\arcsin v+v\sqrt{1-v^2}}{\pi(1+v)^2}=:\Phi(v).
\]
Here $r\mapsto v$ is a bijection of $(0,1)$ onto itself, reversing order.
Writing $A(v)=\arcsin v+v\sqrt{1-v^2}$ gives
\[
A'(v)=2\sqrt{1-v^2},\qquad
\Phi'(v)=\frac{2(\sqrt{1-v^2}-\arcsin v)}{\pi(1+v)^3}.
\]
With $t=\arcsin v$, the numerator changes sign exactly at $t=\tau$.
Thus $v=\sin\tau$ is the unique global maximum. The endpoint limits
are $0$ and $1/8$, and strict monotonicity on either side excludes
boundary maximizers. At this point $\sqrt{1-v^2}=\tau$, so
$A(v)=\tau(1+v)$ and $\Phi(v)=\tau/(\pi(1+v))$.
Since $1+q=2/(1+v)$, this proves \eqref{eq:J} and the claimed value.
Each $D_i$ is now an ordinary integral of a continuous, exactly specified
function over a uniquely specified interval. Its sign and enclosure are
proved arithmetically in Section~\ref{sec:arithmetic}.
\end{proof}
The optimization just made concerns a terminal integral. The lower-bound
proof below uses its identity \eqref{eq:J}; no theorem about an exactly
minimizing finite Supnick cycle is required as an additional premise.

\section{One outer order, its dual and its energy (L3)}\label{sec:dual}
Let $n\ge102$, $k=\lfloor qn\rfloor$, $\alpha=k/n$, $s=1+\alpha$,
$V=\{k/n,(k+1)/n,\ldots,1\}$. The strict bound $q>3/17$ gives
\begin{equation}\label{eq:gridgate}
\alpha>q-1/n>3/17-1/102=1/6,\qquad \alpha\le q<1/5.
\end{equation}
Fix one genuine cyclic order $\sigma$ of the original labels $k,\ldots,n$.
Use normalized coordinates $x=j/n$ in its edges and set
\[
\begin{gathered}
f(x,y)=\sqrt{xy},\qquad r(x)=s-x,\\
W_0=\frac1n\sum_{\{x,y\}\in\sigma} f(x,y),\quad
J_n=\frac1n\sum_{x\in V} f(x,r(x)),\quad e=W_0-J_n.
\end{gathered}
\]
Thus $W_0=S(\sigma)/n^2$. The reflection $r$ permutes $V$ exactly.

\begin{lemma}[Assignment dual]\label{lem:dual}
For $c=s/2$ define
\[
\psi(x)=\frac c2+\frac12\int_c^x\sqrt{\frac{s-t}{t}}\dd t.
\]
Then for $x,y\in[\alpha,1]$,
\begin{equation}\label{eq:dual}
\frac{(x+y-s)^2}{8}\le f(x,y)-\psi(x)-\psi(y)
\le\frac{(x+y-s)^2}{8\alpha}.
\end{equation}
Consequently
\begin{equation}\label{eq:E}
e\ge0,\qquad E:=\frac1n\sum_{\{x,y\}\in\sigma}(x+y-s)^2\le8e.
\end{equation}
\end{lemma}
\begin{proof}
Differentiation shows that $\psi(x)+\psi(s-x)$ and $f(x,s-x)$ have
equal derivatives; at $x=c$ both equal $c$. For
$H(x,y)=f(x,y)-\psi(x)-\psi(y)$ this gives $H(x,s-x)=0$, and
\[
\partial_yH(x,y)=
\frac{x+y-s}{2\sqrt y(\sqrt x+\sqrt{s-y})}.
\]
The denominator lies between $4\alpha$ and $4$. Integrate from $s-x$
to $y$. If $y<s-x$, the integration direction and numerator both
reverse sign. In both cases integrating the linear numerator gives
\eqref{eq:dual}. Reflection and degree two imply
$\sum_{x\in V}f(x,r(x))=2\sum_{x\in V}\psi(x)
=\sum_{\{x,y\}\in\sigma}(\psi(x)+\psi(y))$.
Summing \eqref{eq:dual} proves \eqref{eq:E}.
\end{proof}
The reference $J_n$ need not be the cost of a realizable cycle:
reflection can have two-cycles and a middle fixed point. Only its exact
assignment identity was used. No loop or duplicate neighbor is inserted
into $\sigma$.

Replace each cyclic edge by its two orientations and define
\begin{equation}\label{eq:measure}
\mu=\frac1{2n}\sum_{(x,y)\text{ oriented edge}}\delta_{(x,r(y))},
\qquad \nu=\frac1n\sum_{x\in V}\delta_x.
\end{equation}
Both marginals of $\mu$ are $\nu$, since each vertex occurs twice in
each endpoint position and reflection permutes the grid. Hence
\begin{equation}\label{eq:marginals}
\int[A(x)-A(z)]\dd\mu(x,z)=0,\qquad
\int(x-z)^2\dd\mu(x,z)=E
\end{equation}
for any function $A$ on $V$. The total mass is $(n-k+1)/n$, not
artificially renormalized to one. Every restriction below uses this
same $\sigma$, $\mu$ and $E$.

\section{Signed deletion, seams and finite errors (L4)}\label{sec:deletion}
Consider any $\beta\in(q,23/100]$ and let
$\ell=\lfloor\beta n\rfloor$. Assume in addition
\begin{equation}\label{eq:cutgate}
\ell/n>q,\qquad \ell\le n-2.
\end{equation}
These gates ensure a nonempty deleted set and at least three survivors.
Write $\gamma=\ell/n$ and $B=(\ell-1/2)/n$. A grid vertex is deleted
exactly when $x<B$. Since
$s-2\gamma>7/6-46/100=53/75>1/2$, deleted vertices lie below $s/2$.
Put
\begin{align*}
D_\alpha(x)&=s-x-2\sqrt{x(s-x)},&
D_n&=\frac1n\sum_{j=k}^{\ell-1}D_\alpha(j/n),\\
W_\beta&=\frac1n\sum_{\{x,y\}\in\sigma|\{\ell,\ldots,n\}}f(x,y),&
\Delta&=W_\beta-W_0-D_n,\\
K&=\int |x-z|\ind_{(x-B)(z-B)<0}\dd\mu(x,z).
\end{align*}

\begin{lemma}[Midpoint strip mass]\label{lem:strip}
For every $h>0$, $\nu\{x:|x-B|\le h\}\le4h$.
\end{lemma}
\begin{proof}
No grid point is within $1/(2n)$ of $B$. Below this radius the set is
empty. Otherwise it has at most $2nh+1$ points, of total mass at most
$2h+1/n\le4h$. Truncating the grid at either endpoint only reduces mass.
\end{proof}

\begin{lemma}[Signed first variation]\label{lem:variation}
Let
\[
p(x)=\tfrac12\left(1-\sqrt{x/(s-x)}\right),\quad
P(x)=p(x)\ind_{x<B},\quad
\Lambda=\frac1n\sum_{x<B}p(x)\sum_{y\sim x}(x+y-s).
\]
The neighbors $y\sim x$ are those in the original outer cycle. Then
\begin{equation}\label{eq:43}
\Lambda=2\int P(x)(x-z)\dd\mu,\qquad
|\Delta-\Lambda|\le43E.
\end{equation}
\end{lemma}
\begin{proof}
The first identity follows directly from the oriented mass $1/(2n)$.
Partition deleted vertices into maximal cyclic runs. Rotate the written
cycle to a survivor if necessary. Every removed edge and every replacement
edge is then accounted for exactly once, including the seam.

For an isolated deleted vertex $x$ with surviving neighbors $y,z$ the
change in unnormalized $f$-cost is
$g(x;y,z)=f(y,z)-f(x,y)-f(x,z)$. At $y=z=t=s-x$ it equals
$D_\alpha(x)$, and both first derivatives in $y,z$ equal $p(x)$.
The Hessian entries are
\[
g_{yy}=\frac{\sqrt{x}-\sqrt z}{4y^{3/2}},\quad
g_{zz}=\frac{\sqrt{x}-\sqrt y}{4z^{3/2}},\quad
g_{yz}=\frac1{4\sqrt{yz}}.
\]
On $[\alpha,1]^3$ their absolute values are respectively less than
$4,4,2$: $1/(4\alpha^{3/2})<3\sqrt6/2<4$ and
$1/(4\alpha)<3/2<2$. The symmetric Hessian therefore has operator
norm at most $6$ (each absolute row sum is at most $6$). Taylor's
integral remainder along the segment from $(t,t)$ to $(y,z)$ yields
\[
|g-D_\alpha(x)-p(x)(y-t+z-t)|
\le3((y-t)^2+(z-t)^2).
\]
The segment stays in $[\alpha,1]^2$. Across isolated deletions each
incident edge occurs at most once, so their total normalized error is
at most $3E$.

For runs of length $m\ge2$, let $M_{LL}$ be the total number of outer
edges with both endpoints deleted and $V_{\rm bad}$ the number of
vertices in these runs. Since $m\le2(m-1)$,
$V_{\rm bad}\le2M_{LL}$. A run replaces $m+1$ edges by one edge;
all weights lie in $[0,1]$, so its cost change has absolute value at
most $m+1\le2m$. Since $|D_\alpha(x)|\le2$, its discrepancy from
the $m$ ideal changes is at most $4m$. The total normalized discrepancy
of these runs is at most $8M_{LL}/n$.
For a deleted vertex, $0<p(x)<1/2$ and each defect has absolute value
at most $1-\alpha<1$. The part of $\Lambda$ from bad vertices therefore
has absolute value at most $V_{\rm bad}/n\le2M_{LL}/n$; this includes
both incidences of each $LL$ edge. On each such edge
$|x+y-s|>1/2$, so $M_{LL}/n\le4E$. Combining contributions gives
$3E+(8+2)M_{LL}/n\le43E$.
\end{proof}

\begin{lemma}[Cancellation and signed deletion]\label{lem:54}
Under the same hypotheses,
\begin{equation}\label{eq:54}
|\Lambda|\le11E+2K,\qquad
|W_\beta-W_0-D_n|\le54E+2K.
\end{equation}
\end{lemma}
\begin{proof}
Throughout $[\alpha,1]$, $|p(x)|<1$ and
\[
p'(x)=-\frac{s}{4\sqrt x(s-x)^{3/2}},\qquad
|p'(x)|\le\frac{s}{4\alpha^2}<\frac{54}{5}<11.
\]
The first bound follows from $\sqrt{x/(s-x)}<\sqrt6<3$;
the second uses $s<6/5$ and $\alpha>1/6$.
Let $A(u)=\int_\alpha^u P(v)\dd v$, a continuous primitive with a
piecewise continuous derivative. For grid points $x,z$, integrate
\[
P(x)-P(v)=\ind_{x<B}(p(x)-p(v))
+p(v)(\ind_{x<B}-\ind_{v<B})
\]
along the segment from $z$ to $x$. The smooth term contributes at most
$11|x-z|^2/2$. The jump term vanishes unless the segment crosses $B$
and otherwise contributes at most $|x-z|$. Both orientations give
\[
|P(x)(x-z)-(A(x)-A(z))|
\le\tfrac{11}{2}(x-z)^2+|x-z|\ind_{(x-B)(z-B)<0}.
\]
Neither endpoint equals the midpoint $B$. Equal marginals in
\eqref{eq:marginals} cancel the primitive exactly. Multiply the remaining
bound by two and use \eqref{eq:43} to obtain $11E+2K$, then add $43E$.
Thus the constant is $43+11=54$, with no mesh error in this estimate.
\end{proof}

\begin{lemma}[Uniform floor estimates]\label{lem:floors}
Set $D(\beta)=\int_q^\beta D_q(x)\dd x$.
For every cutoff satisfying \eqref{eq:cutgate},
\begin{equation}\label{eq:floorerrors}
|J_n-\pi\Cterm|\le5/n,\qquad |D_n-D(\beta)|\le19/(2n).
\end{equation}
\end{lemma}
\begin{proof}
Let $F_\alpha(x)=\sqrt{x(1+\alpha-x)}$. All factors in the
comparisons below lie in $[1/6,1]$, and direct differentiation gives
\[
|\partial_xF_\alpha|\le3,\quad
|\partial_\alpha F_\alpha|\le3/2,\quad |F_\alpha|\le1,
\qquad
|\partial_xD_\alpha|\le7,\quad
|\partial_\alpha D_\alpha|\le2,\quad |D_\alpha|\le2.
\]
For example each coordinate derivative of $\sqrt{xt}$ is at most
$\sqrt6/2<3/2$, and
$\partial_\alpha D_\alpha=1-\sqrt{x/(1+\alpha-x)}$.
For a Lipschitz function with constant $M$, the left-rectangle error
on each mesh interval is at most $M/(2n^2)$. There are at most $n$
intervals. The inclusive sum $J_n$ also has its last endpoint term;
its error against $\int_\alpha^1F_\alpha$ is at most $(3/2+1)/n$.
Removing $[\alpha,q]$ costs at most $1/n$ and changing the parameter
from $\alpha$ to $q$ on $[q,1]$ costs at most $3/(2n)$. Their sum is $5/n$.

The sum $D_n$ is exactly a left sum on $[\alpha,\gamma]$; its error
is at most $7/(2n)$. Removing $[\alpha,q]$, changing the parameter on
$[q,\gamma]$ and adding $[\gamma,\beta]$ cost at most $2/n$ each.
Here $\gamma>q$ is an explicit hypothesis. Their sum is $19/(2n)$.
Only $0\le q-\alpha<1/n$ and $0\le\beta-\gamma<1/n$ are used.
No numerical evaluation of a transcendental floor or parity selection occurs.
\end{proof}

\section{One shared crossing budget (L5)}\label{sec:crossings}
\begin{lemma}\label{lem:crossing}
Let $B_1<\cdots<B_m$, $h_i>0$, and let $\mu$ be a finite nonnegative
measure with first marginal $\nu$, finite energy
$E=\int|x-z|^2\dd\mu$ and strip bounds
$\nu\{|x-B_i|\le h_i\}\le4h_i$. Define
$K_i=\int|x-z|\ind_{(x-B_i)(z-B_i)<0}\dd\mu$.
If $h_i+h_{i+1}\le B_{i+1}-B_i$ for every adjacent pair, then
\begin{equation}\label{eq:shared}
\sum_{i=1}^m h_iK_i\le4\sum_{i=1}^m h_i^3+E.
\end{equation}
\end{lemma}
\begin{proof}
Write $d=|x-z|$. At cutoff $i$ split crossings into short ($d\le h_i$)
and long ($d>h_i$). A short crossing has $|x-B_i|\le h_i$, so its
weighted integrand is at most $h_i^2\ind_{|x-B_i|\le h_i}$. Integrating
and summing gives at most $4\sum h_i^3$.
The strictly crossed cutoffs form a consecutive block $p,\ldots,r$.
With no crossing the long contribution is zero; with one it is at most
$h_pd<d^2$ if long. If $r>p$, summing adjacent separation gives
\[
\sum_{i=p}^r h_i
=\sum_{i=p}^{r-1}(h_i+h_{i+1})-\sum_{i=p+1}^{r-1}h_i
\le B_r-B_p-\sum_{i=p+1}^{r-1}h_i\le B_r-B_p<d.
\]
Thus every crossing in this block is long, and their combined weighted
integrand is $d\sum_{i=p}^rh_i\le d^2$. Integration spends $E$ once,
even when the same pair crosses all three thresholds. The equality case
$d=h_i$ was included in the short part; endpoints at a cutoff do not
strictly cross it. No pair is omitted or charged to several energy budgets.
\end{proof}

\section{Finite minimax and transfer to the global problem (L6)}\label{sec:minimax}
Use the fixed three cutoffs of \eqref{eq:parameters}. Let
$\ell_i=\lfloor\beta_i n\rfloor$, $B_i=(\ell_i-1/2)/n$ and
$\sigma_i=\sigma|\{\ell_i,\ldots,n\}$, with $\sigma_0=\sigma$.
Put $W_i=S(\sigma_i)/n^2$ and let $D_{i,n}$ denote $D_n$ at cutoff $i$.
Define the common-order minimax
\[
M_{n,4}=\min_{\sigma\text{ cyclic on }\{\lfloor qn\rfloor,\ldots,n\}}
\max_{0\le i\le3}\Rc(\sigma_i).
\]
Only sufficiently large $n$ are needed; all its cycles then have at least
three positive radii. For positive widths define
\[
H=\sum_i h_i,\quad Q_3=\sum_i h_i^3,\quad L=16+432H,
\quad F_n=\sum_i h_iD_{i,n}-8Q_3.
\]

\begin{proposition}\label{prop:finite}
Suppose
\begin{equation}\label{eq:finitegates}
n\ge102,\quad \ell_1/n>q,\quad
\ell_1<\ell_2<\ell_3\le n-2,\quad
h_i+h_{i+1}\le\frac{\ell_{i+1}-\ell_i}{n}\quad(i=1,2).
\end{equation}
Then
\begin{align}
\Rstar(n)&\ge M_{n,4},\label{eq:globalfinite}\\
\frac{M_{n,4}}{n^2}&\ge\frac{J_n+(F_n)_+/L}{\pi}-\frac1n
\ge\Cterm+\eta(h)-\frac{A(h)}n,\label{eq:finite}\\
A(h)&=1+\frac{5+19H/(2L)}\pi.\nonumber
\end{align}
\end{proposition}
\begin{proof}
All three cutoffs satisfy Section~\ref{sec:deletion}'s hypotheses. With
the same outer measure, Lemmas~\ref{lem:54} and \ref{lem:crossing} give
\begin{align*}
\sum_i h_iW_i
&\ge H(J_n+e)+\sum_i h_iD_{i,n}-54HE-2\sum_i h_iK_i\\
&\ge H(J_n+e)+F_n-(54H+2)E\\
&\ge H(J_n+e)+F_n-Le.
\end{align*}
Here $L=8(54H+2)$ because $E\le8e$; the factor $16$ is not
multiplied by the number of cutoffs. Retaining the outer cost and the
positive weighted mean of the inner costs yields
\[
\max_{0\le i\le3}W_i\ge J_n+
\max\left\{e,\frac{F_n}{H}-\left(\frac LH-1\right)e\right\}.
\]
Since $L>H>0$, if $F_n>0$ the increasing and decreasing branches meet
at $e=F_n/L\ge0$. Below that point the decreasing branch is at least
its meeting value, and above it the increasing branch is at least that
value. If $F_n\le0$, the first branch is always nonnegative and $e=0$
makes the maximum zero. The infimum over all $e\ge0$ is therefore
$(F_n)_+/L$ in every case. This scalar infimum need not be realized by
a tour. The lower sandwich \eqref{eq:sandwich} and a finite minimum
over orders prove the first inequality in \eqref{eq:finite}.

Lemma~\ref{lem:floors} gives $F_n\ge F(h)-19H/(2n)$ and
$J_n\ge\pi\Cterm-5/n$. The positive-part function is increasing and
1-Lipschitz, so $(F_n)_+\ge F(h)_+-19H/(2n)$. This proves the second
inequality even when $F_n$ or $F(h)$ is nonpositive.

For \eqref{eq:globalfinite}, take any full feasible original placement
at radius $R$. Delete the labels below $\lfloor qn\rfloor$ to obtain
one common order $\sigma$. Its three further restrictions are exactly
$\sigma_i$. Lemma~\ref{lem:deletiongeometry} bounds each chain root by
$R$. Thus $M_{n,4}\le R$ for every full feasible placement. Taking the
infimum over $R$ proves the assertion without an attainment assumption.
\end{proof}
All estimates are uniform in the order, include the seam and consecutive
deleted runs, and apply to both cardinality parities. No implication from
chain feasibility to full feasibility was used.

\section{The unique global width optimum (L7)}\label{sec:width}
Set $d=a+b$ and $\ell=16+432d$; these are width constants, not grid
parameters. On the edge where both separation constraints are active set
\begin{align}
N(z)&=F(a-z,z,b-z)=c_0+c_1z-24dz^2+8z^3,\nonumber\\
c_0&=aD_1+bD_3-8(a^3+b^3),\qquad
c_1=D_2-D_1-D_3+24(a^2+b^2),\nonumber\\
L(z)&=\ell-432z,\nonumber\\
G(z)&=N'(z)L(z)+432N(z)\label{eq:G}\\
&=\ell c_1+432c_0-48\ell dz+(24\ell+10368d)z^2-6912z^3.\nonumber
\end{align}
Its exact derivative is
\begin{equation}\label{eq:Gprime}
G'(z)=N''(z)L(z)=48(z-d)L(z)<0\qquad(0\le z\le a),
\end{equation}
since $a<d$ and $L(z)\ge16+432b>0$.

The rational computations proved in Section~\ref{sec:arithmetic} give
$G(0)>0$, $G(a)<0$ and the more precise opposite signs at
\begin{equation}\label{eq:zbracket}
z_-=0.00451910758124826,\qquad z_+=0.00451910758124827.
\end{equation}
Together with \eqref{eq:Gprime}, continuity proves the existence and
uniqueness of $z_*\in(z_-,z_+)$. All coordinates of
$u=(a-z_*,z_*,b-z_*)$ are strictly positive.

\begin{lemma}[Global cubic certificate]\label{lem:KKT}
Let $\rho=F(u)/\mathcal L(u)$ and
$g_i=D_i-24u_i^2-432\rho$. Then $F(u)>0$ and
\[
\lambda_1=g_1>0.00035,\qquad \lambda_2=g_3>0.00059,
\qquad(g_1,g_2,g_3)=(\lambda_1,\lambda_1+\lambda_2,\lambda_2).
\]
For every $h\in\mathcal P$,
\begin{align}\label{eq:cubiccertificate}
F(h)-\rho\mathcal L(h)
={}&\lambda_1(h_1+h_2-a)+\lambda_2(h_2+h_3-b)\nonumber\\
&-8\sum_{i=1}^3(h_i-u_i)^2(h_i+2u_i)\le0,
\end{align}
with strict inequality unless $h=u$. Thus $u$ is the unique global
maximizer of $\eta$ on the entire closed polyhedron.
\end{lemma}
\begin{proof}
The positivity gates are rational consequences of the integral and root
intervals in Section~\ref{sec:arithmetic}. Since $G(z_*)=0$, division
by $L(z_*)$ gives $N'(z_*)+432\rho=0$. On substituting $N'$ this is
$-g_1+g_2-g_3=0$. The gradient of $F-\rho\mathcal L$ at $u$ is
therefore the stated positive combination of the two constraint normals.
For each coordinate the exact identity
\[
h_i^3-u_i^3-3u_i^2(h_i-u_i)=(h_i-u_i)^2(h_i+2u_i)
\]
gives \eqref{eq:cubiccertificate}, since $F(u)-\rho\mathcal L(u)=0$
and the two constraints are equalities at $u$. The linear terms are
nonpositive on $\mathcal P$. Because $h_i\ge0$ and $u_i>0$, the cubic
remainder is strictly negative if any coordinate differs. Hence
$F(h)/\mathcal L(h)<\rho$ for $h\ne u$. If $F(h)\le0$, the
positive-part objective is zero and is also strictly smaller than
$\rho/\pi>0$. This covers zero widths, all faces and vertices, and the
whole nonpositive region without assuming beforehand that an optimum
must lie on the chosen edge. Existence also follows directly from
compactness of $\mathcal P$ and continuity of $\eta$.
\end{proof}

\section{Reproducible rational arithmetic (L2 and L7)}\label{sec:arithmetic}
The included \texttt{check\_width\_optimum.py} uses only Python's standard
library integers and \texttt{Fraction}. It starts from the rational
intervals below, proves that they contain the named exact constants,
and does not import a saved numerical result or an optimizer. To make
the arithmetic reproducible by an independent implementation, we specify
every enclosure formula.

\subsection{Trigonometric gates}
Use the intervals
\begin{align*}
T_-&=0.7390851332151606416553120876738734040134,\\
T_+&=0.7390851332151606416553120876738734040135,\\
q_-&=0.1950200913506069300798071259134151019366,\\
q_+&=0.1950200913506069300798071259134151019367,\\
\pi_-&=3.1415926535897932384626433832795028841971,\\
\pi_+&=3.1415926535897932384626433832795028841972.
\end{align*}
Let
\[
C_m(x)=\sum_{j=0}^{m-1}\frac{(-1)^jx^{2j}}{(2j)!},\quad
S_m(x)=\sum_{j=0}^{m-1}\frac{(-1)^jx^{2j+1}}{(2j+1)!},\quad
A_m(x)=\sum_{j=0}^{m-1}\frac{(-1)^jx^{2j+1}}{2j+1}.
\]
For $0<x<1$ the terms decrease to zero, so an even number of terms
is a lower bound and an odd number is an upper bound. Exact rational
comparisons give
\[
C_{30}(T_-)>T_-,\qquad C_{31}(T_+)<T_+.
\]
Monotonicity of $\cos t-t$ isolates $\tau$. Sine is increasing on
$(0,1)$, so $s_-=S_{30}(T_-)<\sin\tau<S_{31}(T_+)=s_+$.
The decreasing map $s\mapsto(1-s)/(1+s)$ and exact comparisons give
\[
q_-<\frac{1-s_+}{1+s_+}<q<\frac{1-s_-}{1+s_-}<q_+.
\]
In particular $q_->3/17$ and $q_+<1/5$ are checked as rational
inequalities.

For $\pi$ use Machin's identity
$\pi=16\arctan(1/5)-4\arctan(1/239)$. Here is a branch justification:
two tangent doublings give $\tan(4\arctan(1/5))=120/119$,
and the tangent of $4\arctan(1/5)-\arctan(1/239)$ is exactly one.
This angle lies strictly between zero and $4/5<\pi/2$ (the elementary
inequality $\pi>2$ follows, for example, by comparing a semicircle
with its diameter). Positivity follows from
$\arctan(1/5)>1/5-(1/5)^3/3$ and $\arctan(1/239)<1/239$.
The angle is therefore $\pi/4$. Alternating sums give the exact gates
\[
\pi_-<16A_{32}(1/5)-4A_{31}(1/239)
<\pi<16A_{31}(1/5)-4A_{32}(1/239)<\pi_+.
\]

\subsection{Integral certificates}
For each $\beta_i$ use $s_0=1+q_-$,
$v=(1-q_-)/s_0$, $w=(s_0-2\beta_i)/s_0$. These satisfy $0<w<v<1$.
The convergent binomial expansion is
\[
\sqrt{1-y^2}=1-\sum_{j\ge1}c_jy^{2j},\qquad
c_j=\frac{\binom{2j}{j}}{4^j(2j-1)}>0.
\]
Its coefficients decrease, since $c_{j+1}/c_j=(2j-1)/(2j+2)<1$.
Termwise integration is valid uniformly on $[w,v]$, which lies inside
the radius of convergence. Define the rational quantities
\begin{align*}
P&=v-w-\sum_{j=1}^{80}c_j\frac{v^{2j+1}-w^{2j+1}}{2j+1},\\
T&=(v-w)c_{81}\frac{v^{162}}{1-v^2},\qquad
I_-=\frac{s_0^2}{4}(P-T),\quad I_+=\frac{s_0^2}{4}P.
\end{align*}
For every $y\in[w,v]$ the omitted positive tail is at most
$c_{81}v^{162}/(1-v^2)$, hence
\[
I_-\le\int_{q_-}^{\beta_i}\sqrt{x(1+q_--x)}\dd x\le I_+.
\]
Write $D(r,\beta)=\int_r^\beta D_r(x)\dd x$. The endpoint and
integrand bounds are $|D_r(r)|\le2$ and $|\partial_rD_r(x)|\le2$.
Differentiation under the integral therefore gives
$|\partial_rD(r,\beta)|\le4$ on $r\in[q_-,q_+]$.
Thus valid rational endpoints for $D_i$ are
\begin{align*}
d_{i,-}&=s_0(\beta_i-q_-)-(\beta_i^2-q_-^2)/2-2I_+-4(q_+-q_-),\\
d_{i,+}&=s_0(\beta_i-q_-)-(\beta_i^2-q_-^2)/2-2I_-+4(q_+-q_-).
\end{align*}
Before further propagation the script enlarges these endpoints to
30-place rationals: for an interval $[l,r]$ it uses
\[
\left[\frac{\lfloor10^{30}l\rfloor-1}{10^{30}},
\frac{\lceil10^{30}r\rceil+1}{10^{30}}\right].
\]
The extra unit makes the printed enclosures strict; it is an outward
rounding, never a decimal approximation used as an exact input value.

\subsection{Polynomial signs, multipliers and final enclosure}
Interval addition and subtraction use matching and opposite endpoints;
multiplication takes the minimum and maximum of all four products.
Division by a positive interval multiplies by its reversed reciprocal.
Positive powers use monotonic endpoints. These operations remain safe
when a variable appears repeatedly, although dependency can widen them.
Evaluate $G$ at $0,a,z_-,z_+$ using the intervals for $D_i$.
Their signs are respectively $+,-,+,-$ with strict rational margins.
With $Z=[z_-,z_+]$ use the interval vector
\[
U=([a-z_+,a-z_-],\ Z,\ [b-z_+,b-z_-]).
\]
Evaluate $F(U)$, $\rho=F(U)/\mathcal L(U)$ and
$D_i-24U_i^2-432\rho$ outwardly. These certify $F(u)>0$ and the
two positive multipliers in Lemma~\ref{lem:KKT}. The accompanying
\texttt{certificate.json} gives the rational endpoints of these gates,
including the four signs, so they can be checked without interpreting
floating output. The enclosed quantities include
\begin{align*}
0.00136820131190299900276&<D_1<0.00136820131190299900279,\\
0.00196196296152142894680&<D_2<0.00196196296152142894683,\\
0.00241410289623904895464&<D_3<0.00241410289623904895467,\\
0.00000038803240421408705&<\etaw<0.00000038803240421408841.
\end{align*}
Finally enclose $\Cterm=\tau(1+q)/(2\pi)$ by multiplying the positive
intervals for $\tau,1+q$ and dividing by opposite endpoints of $2\pi$.
Add the interval for $\rho/\pi$. Exact comparisons with the two rationals
in \eqref{eq:mainbound} prove the strict enclosure in that equation.
The script \texttt{verify.py} checks those literal target endpoints, the
certificate content and the input/output manifest; a replay that merely
prints a plausible decimal is not the acceptance condition.

\section{The ordered passage to the limit (L8)}\label{sec:limit}
For any fixed positive width vector with strict macroscopic margins
$\delta_i=\beta_{i+1}-\beta_i-h_i-h_{i+1}>0$, choose an integer $N$
such that $N\ge102$, $N(\beta_1-q)>1$, and $N\delta_i\ge1$.
For every $n\ge N$,
\[
\ell_1/n>\beta_1-1/n>q,\qquad
\frac{\ell_{i+1}-\ell_i}{n}>
\beta_{i+1}-\beta_i-1/n\ge h_i+h_{i+1}.
\]
Also $n-\ell_3\ge77n/100\ge2$. The gates \eqref{eq:finitegates}
hold for every such integer $n$. Proposition~\ref{prop:finite} therefore
implies, at fixed $h$,
\begin{equation}\label{eq:fixedlimit}
\liminf_{n\to\infty}\frac{M_{n,4}}{n^2}\ge\Cterm+\eta(h),\qquad
\liminf_{n\to\infty}\frac{\Rstar(n)}{n^2}\ge\Cterm+\eta(h).
\end{equation}

The optimum $u$ saturates both macroscopic constraints and cannot be
substituted blindly into the finite gates. For any fixed $0<t<1$, use
$h=tu$. Its strict margins are $(1-t)a$ and $(1-t)b$. The rational
gates $\beta_1-q>a$ and $1/a>102$ show that
\[
N(t)=\left\lceil\frac1{(1-t)a}\right\rceil
\]
is a sufficient threshold. Apply \eqref{eq:fixedlimit} with this fixed
$t$ and only then let $t\uparrow1$. By continuity $\eta(tu)\to\eta(u)$,
so
\begin{equation}\label{eq:liminf}
\liminf_n M_{n,4}/n^2\ge\Cterm+\etaw,\qquad
\liminf_n\Rstar(n)/n^2\ge\Cterm+\etaw.
\end{equation}
The quantity $M_{n,4}$ depends on the three cutoffs and common order,
not on the auxiliary widths. The same liminf is bounded for every
fixed $t$; no minimum over orders has been exchanged with a limit.

For completeness, unchanged $u$ fails both finite separation gates at
every $n=50000j+1$, $j\ge1$. The exact floors are
\[
(\ell_1,\ell_2,\ell_3)=(10455j,10907j,11500j),
\]
and their normalized differences are respectively $a(n-1)/n<a$ and
$b(n-1)/n<b$. This explains why the order of limits is essential.
It is not a failure of \eqref{eq:liminf}.

\section{Existence of the global limit}\label{sec:existence}
We include only the line comparison and concatenation argument needed
to identify the liminf in \eqref{eq:liminf} with $C_*$. No finite-word
linear programming characterization or upper-bound construction is needed.

For a finite list of nonnegative marks $v_1,\ldots,v_N$, let $b(v)$ be
the least span of line points carrying those marks, with
\begin{equation}\label{eq:line}
|x_i-x_j|\ge\sqrt{v_iv_j}\quad(i\ne j).
\end{equation}
Equal marks still label distinct objects; zero marks may coincide.
For a fixed order $w$, the minimal increasing placement is
\[
x_1=0,\qquad x_i=\max_{j<i}\{x_j+\sqrt{w_iw_j}\}.
\]
It satisfies all constraints. Induction shows that every feasible
placement in this order starting at zero has coordinates at least these.
Thus the minimum span for a fixed order is attained; minimizing over
finitely many orders proves attainment of $b(v)$. Deleting or decreasing
marks cannot increase $b$; scaling marks by $h\ge0$ scales $b$ by $h$.
For $h>0$ the latter follows by scaling positions in both directions,
and for $h=0$ both sides are zero. These assertions also hold for marks
temporarily larger than one.

For marks in $[0,1]$, let $B(v)$ be the infimum circumference of a circle
on which both directed arcs between each pair are at least
$\sqrt{v_iv_j}$. Cutting a feasible circle gives a feasible line
placement. Conversely close an optimal line placement with a gap of
length one: one arc is the line distance and the complementary arc
contains the new gap, at least $1\ge\sqrt{v_iv_j}$. Hence
\begin{equation}\label{eq:bB}
b(v)\le B(v)\le b(v)+1.
\end{equation}
Write $b_n=b(1/n,2/n,\ldots,1)$. From a full feasible original circle
at radius $R$, scale angles by $(R+n)/(2n)$. The lower kernel bound
in Lemma~\ref{lem:sandwich} produces a $B$-circle of circumference
$\pi(R+n)/n$. Conversely a $B$-circle of circumference $b_n+1$ exists
by the preceding gap construction. Scale it to angular circumference
$2\pi$ and set $R=n(b_n+1)/\pi$. Every directed arc is then at least
$2\sqrt{uv}/R$, hence at least $\theta_R(u,v)$ by the upper kernel
bound. The cosine law gives full original feasibility for all pairs.
Taking infima proves for every $n\ge3$
\begin{equation}\label{eq:globalsqueeze}
\frac{b_n}{\pi n}-\frac1n\le\frac{\Rstar(n)}{n^2}
\le\frac{b_n+1}{\pi n}.
\end{equation}

\begin{lemma}\label{lem:existence}
The limit $E_\infty=\lim_n b_n/n$ exists and equals
$\inf_{n\ge2}(b_n+1)/n$. Consequently $C_*=E_\infty/\pi$ exists.
\end{lemma}
\begin{proof}
Fix $n\ge2$. For any $N\ge n$, put $k=\lceil N/n\rceil$ and
concatenate $k$ copies of an optimal $b_n$ arrangement with inter-block
gaps one. All cross-block distances are at least one, so all pair
constraints hold. Its span is $k(b_n+1)-1$, with exactly $k$ copies
of each of the $n$ marks. Retain the $N$ smallest marks. In sorted
mark order their $i$th value is $\lceil i/k\rceil/n$.
Scale all positions and marks by $h=nk/N$, so this value is at least
$i/N$. Decrease the retained marks to $i/N$, assigning each label once
(ties may be broken arbitrarily). Feasibility is preserved. Thus
\[
b_N\le\frac{nk}{N}\,[k(b_n+1)-1].
\]
Since $k/N\to1/n$, division by $N$ gives
$\limsup_N b_N/N\le(b_n+1)/n$ for every fixed $n$.
Equally spaced points with gaps one give $0\le b_N/N\le1$.
Choose $n$ along an unbounded subsequence attaining the liminf; then
$\limsup_N b_N/N\le\liminf_n b_n/n$, proving existence.
The preceding upper bounds and $n\to\infty$ identify the limit with
the displayed infimum. Equation~\eqref{eq:globalsqueeze} now proves
the final assertion for the full geometric problem.
\end{proof}
Combining Lemma~\ref{lem:existence}, \eqref{eq:liminf} and the rational
enclosure proves Theorem~\ref{thm:main}. The concatenation construction
uses actual labeled points and every integer $N$, not only a subsequence
or a coupling with correct marginals.

\section{Verification scope and provenance}\label{sec:scope}
The proof establishes exact theorems and corollaries. The parameter
enclosures are finite rational computations supporting their stated
signs and constants. Bounded tour and measure checks test implementation
invariants and failure modes; they do not prove the universal statements
by enumeration. Neither endpoint of a geometric asymptotic interval is
claimed sharp. The optimized width value is sharp only for the function
$\eta$ on the specified polyhedron at the fixed cutoffs.

The analytic provenance is the immutable scientific snapshot
\begin{center}\small\texttt{3beb8d70c5b3748d370a92855847bdf574e5a14f}.\end{center}
The four source notes are identified in the bibliography. The published
target is\par\smallskip
{\small\path{paper_assets/asymptotic_sequel/ringmin_asymptotic.tex}},\par\smallskip
Proposition 7; the older finite-paper v2 source is not used as the target.
The package includes a dependency map, \path{DEPENDENCIES.md}, recording source sections,
hypotheses and verification limits. Its \texttt{MANIFEST.json} records
input and output SHA-256 hashes, and \texttt{VERIFICATION.md} reports
only checks actually run. These supporting files are included alongside
this source and PDF; no checkout history or old task dossier is required.

\begin{thebibliography}{9}\small
\bibitem{paper} M. Falconi, Ringmin asymptotic sequel,
\href{https://arxiv.org/abs/2609.13630}{arXiv:2609.13630}, Proposition 7.
\bibitem{terminal} M. Falconi, \emph{Optimized induced-terminal-subset
asymptotic lower bound}, Sections 5--7, scientific snapshot above;
\texttt{research/INDUCED\_SUBSET\_ASYMPTOTIC\_LOWER\_BOUND.md}.
\bibitem{stability} M. Falconi, \emph{Quantitative stability for one common
pair of terminal chain tours}, Sections 2--5, 11.1--11.3, same snapshot;
\texttt{research/COMMON\_CHAIN\_QUANTITATIVE\_STABILITY.md}.
\bibitem{crossing} M. Falconi, \emph{A three-level common-chain lower bound
from two shared cutoff budgets}, Sections 7, 9, 11--12, same snapshot;
\texttt{research/THREE\_LEVEL\_COMMON\_CHAIN.md}.
\bibitem{limit} M. Falconi, \emph{The global asymptotic limit and finite-word
variational certificates}, Sections 1--3, same snapshot;
\texttt{research/GLOBAL\_ASYMPTOTIC\_VARIATIONAL\_LIMIT.md}.
\end{thebibliography}
\end{document}
